\chapter{JUSTIFICACIÓN}

El desarrollo del Sistema de Gestión de Turnos para Centros de Salud (San Marcos Carazo) surge como respuesta a la necesidad de mejorar la organización y eficiencia en la atención médica. En muchos centros de salud (San Marcos), el proceso de asignación de turnos se realiza de forma manual, lo que genera largas filas, desorden en la atención y dificultades para priorizar a pacientes con condiciones urgentes.

Esta problemática afecta tanto a los pacientes como al personal médico. Por un lado, los pacientes experimentan tiempos de espera prolongados e incertidumbre sobre cuándo serán atendidos. Por otro lado, el personal de salud de (San Marcos) enfrenta dificultades para gestionar adecuadamente el flujo de pacientes, lo que puede provocar estrés laboral y disminuir la calidad del servicio.

La implementación de este prototipo permite automatizar la asignación de turnos y establecer un sistema de prioridades basado en condiciones como emergencias, adultos mayores y mujeres embarazadas. De esta manera, se garantiza una atención más justa, organizada y eficiente.

Además, este sistema contribuye a la modernización de los centros de salud mediante el uso de herramientas tecnológicas, facilitando la toma de decisiones a través de reportes y mejorando el control del proceso de atención.
